\RequirePackage{plautopatch} % 海外製パッケージとの互換性を保つパッチ（\documentclass より前に書く）
\documentclass[uplatex,dvipdfmx]{jlreq}

% --- 数理系必須パッケージ ---
\usepackage{amsmath,amssymb,amsthm,mathtools,bm,physics}
\usepackage{graphicx}
\usepackage{hyperref}
\usepackage{pxjahyper}

% --- 定理環境 ---
\theoremstyle{definition}
\newtheorem{theorem}{定理}[section]
\newtheorem{definition}[theorem]{定義}
\renewcommand{\proofname}{\textbf{証明}}

\title{LaTeX記述サンプル（見本）}
\author{数理 太郎}
\date{\today}

\begin{document}
\maketitle

\section{はじめに}
これは \texttt{uplatex} と \texttt{jlreq} クラスを用いた記述見本です。
レポートを書く際は \texttt{main.tex} を使用してください。

\section{数式の書き方サンプル}
\begin{theorem}[Eulerの公式]
任意の $\theta \in \mathbb{R}$ に対して次が成り立つ：
\begin{equation}
    e^{i\theta} = \cos\theta + i\sin\theta
\end{equation}
\end{theorem}

\section{参考文献の引用サンプル}
参考文献を引用するには \texttt{\textbackslash cite\{キー\}} コマンドを使います。
引用キー（ID）は \texttt{references.bib} ファイル内で定義されています。

例：LaTeXの必読書といえば『美文書作成入門』\cite{okumura2023}です。
TeXそのものの設計については、開発者自身の著書\cite{knuth84}が詳しいです。

% 参考文献（junsrt: 引用した順に番号を振るスタイル）
\bibliographystyle{junsrt}
\bibliography{references} % references.bib を拡張子なしで指定
\end{document}
